%==============================================================================
% CHAPTER 12: Cellular Charge Trajectories and Disease States
%==============================================================================

\chapter{Cellular Charge Trajectories and Disease States}
\label{chap:charge_trajectories}

\epigraph{%
  Disease is not an invader.\\
  It is the cell forgetting\\
  where it ought to be.%
}{--- the author}

\begin{chapterbox}
The $\Sspace$-entropy space $[0,1]^3$ introduced in
Chapter~\ref{chap:sentropy_space} provides a phase space for cellular state
dynamics. Healthy cells follow a closed attractor $\mathcal{A}_{\mathrm{health}}$
--- a stable limit cycle in $(\Sk, \St, \Se)$ coordinates --- maintained by
the restoring forces of the metabolic hierarchy (Chapter~\ref{chap:metabolic_computer})
and the ion-transport system (Chapter~\ref{chap:charge_current}).
Disease is trajectory deviation: all pathological states correspond to
$d_{\Sspace}(\mathrm{trajectory}, \mathcal{A}_{\mathrm{health}}) > d_{\mathrm{threshold}}$.
Three canonical disease classes arise from three orthogonal deviations:
cancer (Warburg effect) as a $\St$ deviation (temporal entropy increase);
neurodegeneration as a $\Sk$ deviation (knowledge entropy increase from ion
pump failure); and aging as a $\Se$ deviation (evolutionary entropy increase
from nuclear pore complex degradation). Each class has a mechanistic
partition-theoretic explanation, a quantitative signature in $\Sspace$, and
a predicted clinical biomarker profile that precedes symptoms.
\end{chapterbox}

%------------------------------------------------------------------------------
\section{The Health Attractor in \texorpdfstring{$\Sspace$}{S}-Space}
\label{sec:health_attractor}
%------------------------------------------------------------------------------

Chapter~\ref{chap:sentropy_space} defined the $\Sspace$-entropy coordinate
$\Scoord = (\Sk, \St, \Se) \in [0,1]^3$ for a cell, where $\Sk$ is knowledge
entropy (ionic and informational order), $\St$ is temporal entropy (metabolic
timing regularity), and $\Se$ is evolutionary entropy (epigenetic memory
fidelity). Each coordinate is a normalised partition depth ratio, bounded by
the Bounded Phase Space Law.

\begin{definition}[Health Attractor]
\label{def:health_attractor}
The \emph{health attractor} $\mathcal{A}_{\mathrm{health}} \subset \Sspace$ is the
set of $\Sspace$-coordinates corresponding to normal cellular function across
all cell types:
\begin{equation}
  \mathcal{A}_{\mathrm{health}} = \bigl\{ \Scoord \in [0,1]^3
    : \text{all metabolic, ionic, and epigenetic homeostatic conditions hold} \bigr\}.
  \label{eq:health_attractor_def}
\end{equation}
For a given cell type $c$, $\mathcal{A}_{\mathrm{health}}^{(c)}$ is a
subset of $\mathcal{A}_{\mathrm{health}}$ specific to that cell's function.
\end{definition}

The attractor is not a point but a closed orbit: cells oscillate through
metabolic states (glycolytic oscillations, circadian rhythms, cell-cycle
checkpoints) while remaining on $\mathcal{A}_{\mathrm{health}}^{(c)}$.

\begin{theorem}[Stability of the Health Attractor]
\label{thm:attractor_stability}
The health attractor $\mathcal{A}_{\mathrm{health}}$ is a stable limit cycle
in $\Sspace$ under normal physiological conditions.
The restoration timescale for a perturbation $\delta\Scoord$ from the attractor
is
\begin{equation}
  \tau_{\mathrm{restore}} = \frac{1}{\Gamma\,|\nabla\Depth|},
  \label{eq:restoration_timescale}
\end{equation}
where $\Gamma$ is the metabolic flux rate and $|\nabla\Depth|$ is the partition
depth gradient magnitude at the perturbed point in $\Sspace$.
\end{theorem}

\begin{proof}
The metabolic hierarchy (Chapter~\ref{chap:metabolic_computer}) provides
restoring forces in all three $\Sspace$-coordinates. For $\Sk$: ion pump
activity (Chapter~\ref{chap:charge_current}) restores ionic partition depth
differences when $\Sk$ deviates. For $\St$: glycolytic and circadian
oscillators are self-correcting limit cycles (Chapter~\ref{chap:metabolic_clock});
perturbations to metabolic timing decay on the oscillation period timescale.
For $\Se$: epigenetic maintenance machinery (DNA methyltransferases, histone
modifiers, nuclear pore complexes) continuously repairs epigenetic marks.

In the linearised approximation near $\mathcal{A}_{\mathrm{health}}$, the
restoring dynamics take the form
\begin{equation}
  \frac{d(\delta\Scoord)}{dt} = -\Lambda\,\delta\Scoord,
  \label{eq:restoring_dynamics}
\end{equation}
where $\Lambda$ is the $3\times 3$ restoration matrix with positive eigenvalues
$\lambda_k,\lambda_t,\lambda_e > 0$. Each eigenvalue is
$\lambda_i = \Gamma_i\,|\nabla_i\Depth|$, the product of the metabolic flux
rate along coordinate $i$ and the partition depth gradient. The solution
$\delta\Scoord(t) = e^{-\Lambda t}\delta\Scoord(0)$ decays exponentially,
giving $\tau_{\mathrm{restore},i} = 1/\lambda_i = 1/(\Gamma_i|\nabla_i\Depth|)$.
The overall restoration timescale is dominated by the slowest eigenvalue,
$\tau_{\mathrm{restore}} = \max_i\tau_{\mathrm{restore},i}$. \qed
\end{proof}

\begin{remark}
Typical restoration timescales: $\tau_{\mathrm{restore},k} \sim$ minutes
(ionic homeostasis, Na$^+$/K$^+$-ATPase timescale); $\tau_{\mathrm{restore},t}
\sim$ hours (metabolic oscillation re-entrainment); $\tau_{\mathrm{restore},e}
\sim$ days to weeks (epigenetic remodelling). Disease represents failure of
these restorative mechanisms at one or more coordinates.
\end{remark}

%------------------------------------------------------------------------------
\section{Disease as Trajectory Deviation}
\label{sec:disease_deviation}
%------------------------------------------------------------------------------

\begin{definition}[$\Sspace$-Distance]
\label{def:sspace_distance}
The $\Sspace$-distance from a cellular trajectory $\gamma(t)$ to the health
attractor is
\begin{equation}
  d_{\Sspace}\!\left(\gamma, \mathcal{A}_{\mathrm{health}}\right)
    = \sup_{t \geq 0}\,\min_{\mathbf{a}\in\mathcal{A}_{\mathrm{health}}}
      \|\gamma(t) - \mathbf{a}\|,
  \label{eq:sspace_distance}
\end{equation}
where $\|\cdot\|$ is the Euclidean norm on $[0,1]^3$.
\end{definition}

\begin{theorem}[Disease as Trajectory Deviation]
\label{thm:disease_deviation}
A cell is in a disease state if and only if
\begin{equation}
  d_{\Sspace}\!\left(\gamma, \mathcal{A}_{\mathrm{health}}\right) > d_{\mathrm{threshold}},
  \label{eq:disease_condition}
\end{equation}
where $d_{\mathrm{threshold}}$ is the noise tolerance of the cellular
homeostatic machinery, determined by the amplitude of normal metabolic
fluctuations projected onto $\Sspace$.
\end{theorem}

\begin{proof}
By Theorem~\ref{thm:attractor_stability}, perturbations with
$\|\delta\Scoord\| < d_{\mathrm{threshold}}$ decay exponentially to
$\mathcal{A}_{\mathrm{health}}$. For perturbations with
$\|\delta\Scoord\| > d_{\mathrm{threshold}}$, the restoring force
$-\Lambda\,\delta\Scoord$ is insufficient to overcome the sustained
perturbation (metabolic rewiring, gene mutation, pump failure), and the
trajectory diverges from $\mathcal{A}_{\mathrm{health}}$. The threshold
$d_{\mathrm{threshold}}$ is set by the smallest eigenvalue of $\Lambda$
times the noise amplitude in $\Sspace$. \qed
\end{proof}

\begin{theorem}[Universal Disease Classification]
\label{thm:disease_classification}
Every disease state corresponds to a sustained deviation in one or more
$\Sspace$-entropy coordinates:
\begin{itemize}
  \item $\Delta\Sk > 0$ (knowledge entropy increase): loss of cellular
        information about ionic and molecular state
        $\to$ neurodegeneration, genetic diseases.
  \item $\Delta\St > 0$ (temporal entropy increase): disruption of metabolic
        timing and oscillation coherence
        $\to$ metabolic diseases, cancer.
  \item $\Delta\Se > 0$ (evolutionary entropy increase): loss of epigenetic
        memory
        $\to$ aging, epigenetic diseases.
  \item $\Delta\Sk, \Delta\St, \Delta\Se > 0$ (all three elevated):
        cancer at late stage.
\end{itemize}
\end{theorem}

The classification is exhaustive because $\Sspace = [0,1]^3$ is the complete
phase space for cellular state: every perturbation from normal function projects
onto at least one of the three axes.

%------------------------------------------------------------------------------
\section{Cancer: The Warburg Effect as a \texorpdfstring{$\St$}{St} Deviation}
\label{sec:cancer_warburg}
%------------------------------------------------------------------------------

\subsection{Metabolic Timing Disruption}
\label{ssec:warburg_timing}

Cancer cells exhibit the Warburg effect: aerobic glycolysis (lactate production
even in the presence of oxygen) \citep{Warburg1956}. Chapter~\ref{chap:metabolic_computer}
identified this as $L3$--$L4$ decoupling, reducing hierarchical depth to
$D \approx 2$. In $\Sspace$-coordinates, this is a $\St$ deviation.

\begin{proposition}[Cancer as $\St$ Deviation]
\label{prop:cancer_st}
The Warburg phenotype corresponds to $\Delta\St > 0$:
the temporal entropy increases because metabolic oscillation coherence is
disrupted when aerobic glycolysis (L2, short timescale
$\tau_{L2} \sim 1$--$10$\,s) replaces oxidative phosphorylation (L4,
timescale $\tau_{L4} \sim 10$--$100$\,s). The multi-timescale coordination
that defines normal metabolic rhythm (Chapter~\ref{chap:metabolic_computer})
collapses to a single short-timescale process.
\end{proposition}

\begin{proof}
(Sketch.) Normal metabolism coordinates five timescale decades
($\tau_{L1}$ through $\tau_{L5}$), contributing $I_{\mathrm{total}} = 7.29$\,bits
of inter-level information compression. With $D = 2$ (Warburg), only
$I_{\mathrm{total}} \approx 2.9$\,bits of coordination remain
(Table~\ref{tab:hierarchical_depth}, Chapter~\ref{chap:metabolic_computer}).
The loss of $7.29 - 2.9 = 4.39$\,bits of timing information is registered as
$\Delta\St > 0$ in $\Sspace$. \qed
\end{proof}

\subsection{pH as the Partition Depth Carrier of the Warburg Effect}
\label{ssec:warburg_ph}

The aerobic glycolysis of cancer cells produces lactic acid, lowering
intracellular and extracellular pH. Tumour interstitial pH is consistently
6.7--7.0, versus normal tissue pH 7.2--7.4: a shift of $\Delta\mathrm{pH} = 0.2$--$0.5$
units \citep{Gatenby2007}. In partition theory, this $\Delta\mathrm{pH}$ is a
change in the proton partition depth.

From Theorem~\ref{thm:nernst_partition} (Chapter~\ref{chap:charge_current}):
\begin{equation}
  \Delta V_{\mathrm{pH}} = \frac{\kB T}{q}\,\Delta\Depth_{\ce{H+}}\,\ln b
    = \frac{2.303\,\kB T}{q}\,\Delta\mathrm{pH}
    \approx 12\text{--}30\ \mathrm{mV}
\end{equation}
per unit pH change. The partition depth change per proton associated with a
$\Delta\mathrm{pH} = 0.3$ shift is
\begin{equation}
  \Delta\Depth_{\ce{H+}}
    = \frac{\Delta\mathrm{pH}}{\log_{10}(b)}
    = \frac{0.3}{\log_{10}(e)}
    \approx 0.69\ \text{(base-$e$ units)}.
  \label{eq:warburg_depth_change}
\end{equation}

This partition depth reduction biases proton-coupled reactions toward glycolytic
rather than oxidative pathways. Specifically, pyruvate dehydrogenase (PDH,
the L2$\to$L3 gateway) is a proton-coupled enzyme: its activity decreases as
intracellular $[\ce{H+}]$ increases. The Warburg effect is therefore
\emph{self-reinforcing}: glycolysis acidifies the cytoplasm, further inhibiting
PDH, further reducing L3 flux, further increasing $\Delta\St$.

\begin{keyresult}
The partition-theoretic interpretation of the Warburg effect:
tumour acidification ($\Delta\mathrm{pH} \approx 0.3$\,units) changes the proton
partition depth by $\Delta\Depth_{\ce{H+}} \approx 0.69$ (base-$e$ units),
biasing all proton-coupled reactions toward glycolysis and creating a positive
feedback loop that locks the cell into the $D = 2$ (Warburg) metabolic attractor.
The $\St$ deviation that results is measurable as loss of glycolytic oscillation
coherence and can be detected from metabolomics before malignant transformation
is histologically apparent.
\end{keyresult}

\subsection{All Three Coordinates Deviate in Advanced Cancer}
\label{ssec:cancer_all_three}

Late-stage cancer exhibits deviations in all three $\Sspace$-coordinates:
\begin{itemize}
  \item $\Delta\St > 0$: Warburg metabolic timing disruption (primary).
  \item $\Delta\Sk > 0$: Genomic instability increases knowledge entropy
        (the cell loses reliable information about its own sequence state
        through \ce{DNA} damage accumulation).
  \item $\Delta\Se > 0$: Epigenetic dysregulation (cancer-associated
        hypomethylation and hypermethylation patterns) increases evolutionary
        entropy.
\end{itemize}
The $\Sspace$-distance from $\mathcal{A}_{\mathrm{health}}$ grows through the
three deviations additively:
\begin{equation}
  d_{\Sspace}^{\mathrm{cancer}}
    = \sqrt{(\Delta\Sk)^2 + (\Delta\St)^2 + (\Delta\Se)^2},
\end{equation}
with $\Delta\St$ dominant at early stage and all three becoming significant at
metastatic stage.

%------------------------------------------------------------------------------
\section{Neurodegeneration: Ion Pump Failure as a \texorpdfstring{$\Sk$}{Sk} Deviation}
\label{sec:neurodegeneration}
%------------------------------------------------------------------------------

\subsection{The Na\texorpdfstring{$^+$}{+}/K\texorpdfstring{$^+$}{+}-ATPase and Knowledge Entropy}
\label{ssec:naka_knowledge}

The Na$^+$/K$^+$-ATPase maintains the ionic partition depth differences
$\Delta\Depth_{\ce{Na+}}$ and $\Delta\Depth_{\ce{K+}}$ (Chapter~\ref{chap:charge_current},
Section~\ref{sec:naka_atpase}). These differences encode the cell's
\emph{knowledge of its ionic state}: a well-maintained gradient represents
low $\Sk$ (the cell knows precisely where its ions are). Pump dysfunction
raises $[\ce{Na+}]_{\mathrm{in}}$ and lowers $[\ce{K+}]_{\mathrm{in}}$,
collapsing $\Delta\Depth_{\ce{Na+}}$ and $\Delta\Depth_{\ce{K+}}$ and
increasing $\Sk$.

\begin{proposition}[Ion Pump Failure as $\Sk$ Deviation]
\label{prop:pump_failure_sk}
Na$^+$/K$^+$-ATPase dysfunction produces $\Delta\Sk > 0$: the knowledge
entropy increases because the cell loses accurate information about its
ionic partition depth state. The resulting membrane depolarisation
$\Delta V_m = V_m^{\mathrm{fail}} - V_m^{\mathrm{healthy}} > 0$ is
a direct voltage readout of the $\Sk$ increase:
\begin{equation}
  \Delta\Sk \propto \Delta V_m / (V_{\mathrm{Nernst}} - V_m^{\mathrm{healthy}}).
  \label{eq:sk_depolarization}
\end{equation}
\end{proposition}

\subsection{The Neurodegenerative Cascade}
\label{ssec:neurodeg_cascade}

Ion pump failure initiates a cascade of partition topology changes:
\begin{enumerate}
  \item \textbf{Membrane depolarisation.} Loss of $\Delta\Depth_{\ce{Na+}}$
        and $\Delta\Depth_{\ce{K+}}$ raises $V_m$ from $-70$\,mV toward
        $0$\,mV. At intermediate depolarisation ($V_m \approx -40$\,mV),
        voltage-gated \ce{Ca^{2+}} channels open.

  \item \textbf{Calcium influx.} The enormous \ce{Ca^{2+}} partition depth
        difference ($\Delta\Depth_{\ce{Ca^{2+}}} \approx 9.5$, from
        Table~\ref{tab:nernst_potentials}, Chapter~\ref{chap:charge_current})
        drives a massive \ce{Ca^{2+}} influx.
        $[\ce{Ca^{2+}}]_i$ rises from $100$\,nM toward $\mu$M levels.

  \item \textbf{Mitochondrial uncoupling.} Elevated $[\ce{Ca^{2+}}]_i$
        opens mitochondrial permeability transition pores (mPTP), dissipating
        the proton partition depth gradient $\Delta\Depth_{\ce{H+}}$
        (Chapter~\ref{chap:charge_current}, Theorem~\ref{thm:atp_synthesis}).
        ATP synthesis collapses.

  \item \textbf{ATP depletion.} ATP loss further impairs the Na$^+$/K$^+$-ATPase
        (which requires ATP), completing a positive feedback loop:
        pump failure $\to$ depolarisation $\to$ \ce{Ca^{2+}} influx
        $\to$ mitochondrial uncoupling $\to$ ATP depletion
        $\to$ pump failure.

  \item \textbf{Apoptosis.} The cascade activates apoptotic pathways
        (cytochrome $c$ release, caspase activation) through
        mitochondrial outer membrane permeabilisation.
\end{enumerate}

In $\Sspace$-coordinates, this cascade is a rapid $\Sk$ deviation
($\Delta\Sk \gg d_{\mathrm{threshold}}$) that crosses all three axes as it
progresses: apoptosis elevates all three entropy coordinates simultaneously,
driving the cellular trajectory outside all bounds of $\mathcal{A}_{\mathrm{health}}$.

\subsection{Alzheimer's Disease}
\label{ssec:alzheimer}

Amyloid-$\beta$ (A$\beta$) oligomers --- the toxic species in Alzheimer's
disease --- interact directly with the Na$^+$/K$^+$-ATPase \citep{Mark1995}.
A$\beta$ reduces pump activity through two partition-theoretic mechanisms:
\begin{enumerate}[label=(\roman*)]
  \item \textbf{Direct inhibition:} A$\beta$ inserts into the plasma membrane
        and disrupts the partition topology of the pump's transmembrane
        segments, increasing $\Depth_{\mathrm{pump}}$ and reducing its
        rate (Eq.~\eqref{eq:gating_rate}).
  \item \textbf{Oxidative stress:} A$\beta$ generates reactive oxygen species
        (ROS) that oxidise the pump's catalytic subunit, increasing
        $\Depth_{\mathrm{pump}}$ through protein conformational damage.
\end{enumerate}

The result is a sustained $\Sk$ elevation:
$\Delta\Sk^{\mathrm{AD}} \propto -\ln(k_{\mathrm{pump}}^{\mathrm{A}\beta}/k_{\mathrm{pump}}^{\mathrm{ctrl}})$,
where $k_{\mathrm{pump}}$ is the pump rate. Synaptic failure follows because
dendritic spines, which require maintained ionic gradients for long-term
potentiation, lose their partition depth precision: the synapse can no longer
perform precise categorical measurements (Chapter~\ref{chap:hydrogen_bonds}).

\begin{keyresult}
In Alzheimer's disease, amyloid-$\beta$ inhibits the Na$^+$/K$^+$-ATPase,
elevating $\Sk$. This elevation precedes amyloid plaque deposition and
clinical cognitive symptoms by years. Measuring $\Sk$ from neuronal ionic
state (accessible via potassium MRI or blood biomarkers of Na$^+$/K$^+$-ATPase
activity) should provide earlier disease prediction than current amyloid-PET
or CSF biomarkers.
\end{keyresult}

\subsection{Amyotrophic Lateral Sclerosis}
\label{ssec:als}

In ALS, TDP-43 aggregation disrupts RNA processing in motor neurons. In
partition terms, TDP-43 is a partition gate for RNA: it controls the
nuclear-cytoplasmic traffic of specific mRNA species. TDP-43 aggregation
impairs this gate ($\Depth_{\mathrm{TDP43}} \to \infty$ for the affected
transcripts), reducing the export of mRNAs encoding ion channel subunits and
motor neuron maintenance factors.

The partition-theoretic mechanism differs from Alzheimer's in the direction
of the $\Sk$ deviation:
\begin{itemize}
  \item In Alzheimer's: $\Sk^{\mathrm{ionic}}$ increases (cellular knowledge
        of ionic state is lost).
  \item In ALS: $\Sk^{\mathrm{transcriptomic}}$ decreases in the affected
        compartment --- the motor neuron loses access to the transcripts it
        needs, reducing the \emph{usable} partition depth of its
        knowledge state --- while $\Sk^{\mathrm{ionic}}$ increases
        secondarily from channel dysregulation.
\end{itemize}
Both diseases cause the cellular trajectory to exit $\mathcal{A}_{\mathrm{health}}$,
but through distinct paths in $\Sspace$: Alzheimer's moves primarily along
$\hat{e}_k$ (ionic knowledge entropy), while ALS moves along a mixed
$\hat{e}_k$-$\hat{e}_e$ direction (ionic + transcriptomic entropy).

%------------------------------------------------------------------------------
\section{Aging: Nuclear Pore Complex Degradation as a \texorpdfstring{$\Se$}{Se} Deviation}
\label{sec:aging_npc}
%------------------------------------------------------------------------------

\subsection{The Nuclear Pore Complex as a Partition Gate}
\label{ssec:npc_partition}

The nuclear pore complex (NPC) is a $\sim$120\,MDa molecular machine
spanning the nuclear envelope and regulating all transport between nucleus
(the partition containing the genome, epigenome, and transcription machinery)
and cytoplasm (the partition containing ribosomes and metabolic machinery).
In partition terms, the NPC is the principal partition gate between the
$\Se$ (evolutionary entropy) compartment and the rest of the cell.

An intact NPC enforces strict partition selectivity:
$\Depth_{\mathrm{NPC}}(\text{cargo}) \ll 1$ for nuclear import/export
signal-bearing proteins, and $\Depth_{\mathrm{NPC}}(\text{non-specific}) \gg 1$
for non-target molecules. This selectivity is the physical basis of
epigenetic memory: the genome is shielded from cytoplasmic noise by
the NPC's high partition depth for non-specific molecules.

\subsection{NPC Degradation with Age}
\label{ssec:npc_degradation}

NPC scaffold proteins (especially Nup88, Nup98, and Nup153) are among the
longest-lived proteins in post-mitotic cells, with half-lives exceeding the
lifespan of the organism in neurons \citep{Savas2012}. They accumulate
oxidative damage over decades without replacement. As a result, NPC selectivity
decreases with age: the partition depth $\Depth_{\mathrm{NPC}}$ for non-specific
molecules decreases, allowing nuclear-cytoplasmic leakage.

\begin{definition}[NPC Degradation Rate]
\label{def:npc_degradation}
The rate of NPC-mediated partition selectivity loss follows
\begin{equation}
  k_{\mathrm{NPC}}(t) = k_0\,e^{\alpha t},
  \label{eq:npc_degradation}
\end{equation}
where $k_0$ is the basal leakage rate at $t = 0$ (young cell), $\alpha$ is
the aging rate constant (organism- and cell-type-specific), and $t$ is
chronological age. The selectivity of the NPC decreases as $e^{-\alpha t}$:
nuclear-cytoplasmic mixing increases exponentially with age.
\end{definition}

\begin{proposition}[Aging as $\Se$ Deviation]
\label{prop:aging_se}
NPC degradation causes $\Delta\Se > 0$: evolutionary entropy increases because
the nuclear partition boundary becomes leaky, allowing non-specific molecular
mixing between nucleus and cytoplasm. This mixing:
\begin{enumerate}[label=(\roman*)]
  \item Allows cytoplasmic factors to access and inappropriately modify
        the epigenome (widespread epigenetic drift with age).
  \item Allows nuclear factors (transcription factors, chromatin regulators)
        to accumulate in the cytoplasm, reducing their nuclear concentration
        and impairing genome-state maintenance.
  \item Leads to loss of heterochromatin (the high-$\Depth$ genomic
        compartment), as heterochromatin-maintaining factors dilute from
        the nucleus.
\end{enumerate}
\end{proposition}

\begin{proof}
The $\Se$ coordinate measures the fidelity of epigenetic memory transmission:
$\Se = H(\text{epigenetic state at } t) / H_{\max}$, where $H$ is the
Shannon entropy of the epigenetic mark distribution. An intact NPC enforces
a low-$\Se$ nuclear environment by excluding epigenetically disruptive
cytoplasmic species. As NPC selectivity degrades (Eq.~\eqref{eq:npc_degradation}),
non-specific nuclear-cytoplasmic mixing increases at rate $k_{\mathrm{NPC}}(t)$,
stochastically perturbing epigenetic marks. The resulting entropy increase is
\begin{equation}
  \frac{d(\Delta\Se)}{dt} = \kappa\,k_{\mathrm{NPC}}(t)
    = \kappa\,k_0\,e^{\alpha t},
  \label{eq:se_increase_rate}
\end{equation}
where $\kappa$ is the epigenetic sensitivity coefficient (the number of
epigenetic marks disturbed per unit of NPC leakage). Integrating:
\begin{equation}
  \Delta\Se(t) = \frac{\kappa\,k_0}{\alpha}\bigl(e^{\alpha t} - 1\bigr),
  \label{eq:se_increase}
\end{equation}
which is the exponentially accelerating epigenetic drift observed in aging
tissues \citep{Hannum2013,Horvath2013}. \qed
\end{proof}

\begin{remark}
Equation~\eqref{eq:se_increase} predicts exponentially accelerating epigenetic
drift, consistent with DNA methylation clock data \citep{Horvath2013}: epigenetic
age acceleration increases faster than linearly with chronological age.
The aging rate constant $\alpha$ is organism-specific and may be modifiable
by interventions that slow NPC protein turnover damage (caloric restriction,
reduced ROS production).
\end{remark}

\subsection{Mechanistic Consequences of $\Se$ Elevation}
\label{ssec:se_consequences}

Three mechanistic consequences of $\Se$ elevation (NPC degradation) are
experimentally established:

\begin{enumerate}
  \item \textbf{Widespread epigenetic drift.} Methylation clock studies
        \citep{Horvath2013} document the progressive divergence of epigenetic
        state from the young-cell baseline. The divergence rate increases with
        age (Eq.~\eqref{eq:se_increase}), consistent with exponential NPC
        degradation.

  \item \textbf{Heterochromatin loss.} Heterochromatin --- the high-$\Depth$
        genomic compartment that silences transposons and maintains cell
        identity --- progressively decondenses with age \citep{Sedivy2013}.
        The partition depth $\Depth_{\mathrm{HC}}$ decreases, allowing
        transposon activation and aberrant transcription.

  \item \textbf{Increased nuclear-cytoplasmic mixing.} Cytoplasmic proteins
        (including histones, which are produced in the cytoplasm and imported
        to the nucleus) show age-dependent redistribution. Partition depth
        analysis predicts that the nuclear-cytoplasmic ratio of any
        partition-regulated protein $X$ changes as:
        \begin{equation}
          \frac{[X]_{\mathrm{nuc}}}{[X]_{\mathrm{cyt}}}(t)
            = \frac{[X]_{\mathrm{nuc}}}{[X]_{\mathrm{cyt}}}(0)
              \cdot e^{-\Depth_{\mathrm{NPC}}^X(t)}
              = \frac{[X]_{\mathrm{nuc}}}{[X]_{\mathrm{cyt}}}(0)
              \cdot e^{-\Depth_0^X e^{-\alpha t}},
          \label{eq:nc_mixing}
        \end{equation}
        where $\Depth_0^X$ is the young-cell partition depth for protein $X$.
\end{enumerate}

%------------------------------------------------------------------------------
\section{\texorpdfstring{$\Sspace$}{S}-Space Diagnostics and Predictive Medicine}
\label{sec:sspace_diagnostics}
%------------------------------------------------------------------------------

The $\Sspace$-trajectory framework makes a strong predictive claim: cellular
$\Sspace$-deviation precedes clinical symptoms because the attractor divergence
begins before structural or functional damage is clinically detectable.

\begin{theorem}[$\Sspace$-Deviation Precedes Symptoms]
\label{thm:sspace_predictive}
For all three canonical disease classes, the $\Sspace$-deviation
$d_{\Sspace}(\gamma, \mathcal{A}_{\mathrm{health}})$ exceeds $d_{\mathrm{threshold}}$
before clinical symptoms manifest, with a lead time of:
\begin{itemize}
  \item Cancer: years (Warburg metabolic shift precedes malignant transformation).
  \item Neurodegeneration: years (ionic homeostasis disruption precedes
        neuronal death).
  \item Aging: decades (NPC degradation is measurable from epigenetic
        clock acceleration).
\end{itemize}
\end{theorem}

\begin{proof}
(For cancer.) The Warburg effect (L3--L4 decoupling, $\Delta\St > 0$) is
detectable in pre-malignant cells from $^{13}$C-glucose metabolomics and
glycolytic oscillation spectroscopy before histological atypia. The
$\St$ deviation accumulates as metabolic rewiring precedes genomic instability
($\Delta\Sk$ elevation). Since malignant transformation requires additional
genomic events, the metabolic ($\St$) signature must precede the genomic
($\Sk$) signature. Clinically, the Warburg metabolic signature is present in
field carcinogenesis --- the metabolically altered tissue surrounding visible
tumours --- demonstrating that $\St$ elevation precedes tumour formation.

(For neurodegeneration.) In Alzheimer's, synaptic dysfunction (detectable
by $^{18}$F-FDG-PET showing reduced glucose uptake in vulnerable regions)
precedes amyloid plaque accumulation by years \citep{Jack2013}. The ionic
homeostasis failure ($\Delta\Sk$ elevation) drives the synaptic glucose
uptake reduction: depolarised neurons have reduced metabolic efficiency
(Chapter~\ref{chap:metabolic_computer}, Section~\ref{sec:hierarchical_depth}).

(For aging.) DNA methylation clocks \citep{Horvath2013} measure epigenetic
age (a proxy for $\Se$) with single-year precision. Epigenetic age
acceleration --- the rate at which $\Se$ increases beyond the chronological
age trajectory --- is predictive of mortality decades in advance, well before
clinical disease. \qed
\end{proof}

\subsection{Operationalising \texorpdfstring{$\Sspace$}{S}-Space Measurement}
\label{ssec:sspace_measurement}

The three $\Sspace$-coordinates are measurable from current biological assays:

\begin{table}[htbp]
\centering
\caption{Measurable proxies for $\Sspace$-entropy coordinates, their
  correspondence to partition-theoretic quantities, and current clinical
  feasibility.}
\label{tab:sspace_measurement}
\begin{tabular}{llll}
\toprule
Coordinate & Partition quantity & Measurable proxy & Status \\
\midrule
$\Sk$ & Ionic partition depth differences & Serum $[\ce{K+}]$, $[\ce{Na+}]$;
         Na$^+$/K$^+$-ATPase activity & Routine \\
      & & Potassium MRI (brain) & Research \\
$\St$ & Metabolic oscillation coherence & $^{13}$C-glucose metabolomics;
         glycolytic oscillation spectroscopy & Research \\
      & Hierarchical depth $D$ & FDG-PET metabolic rate &
        Routine (oncology) \\
$\Se$ & Epigenetic mark fidelity & DNA methylation clock (blood) &
        Validated \citep{Horvath2013} \\
      & NPC integrity & Nup88 nuclear-cytoplasmic ratio & Research \\
\bottomrule
\end{tabular}
\end{table}

A composite $\Sspace$-diagnostic score could be computed from these proxies
as a single early-disease index. The framework predicts that such a score
would outperform current disease-specific biomarkers in longitudinal prediction
of cancer, neurodegeneration, and accelerated aging.

%------------------------------------------------------------------------------
\section{Summary}
\label{sec:trajectories_summary}
%------------------------------------------------------------------------------

\begin{chapterbox}
The $\Sspace$-entropy space $[0,1]^3$ provides a complete phase space for
cellular state dynamics. Healthy cells follow the stable limit cycle
$\mathcal{A}_{\mathrm{health}}$ maintained by ionic, metabolic, and epigenetic
homeostatic machinery, with restoration timescale
$\tau_{\mathrm{restore}} = 1/(\Gamma|\nabla\Depth|)$
(Theorem~\ref{thm:attractor_stability}). Disease is trajectory deviation:
$d_{\Sspace}(\gamma, \mathcal{A}_{\mathrm{health}}) > d_{\mathrm{threshold}}$
(Theorem~\ref{thm:disease_deviation}). Three canonical classes arise from
orthogonal deviations (Theorem~\ref{thm:disease_classification}):

\medskip
\noindent\textbf{Cancer} ($\Delta\St > 0$): The Warburg effect decouples the
metabolic hierarchy (L3--L4 decoupling, $D \to 2$), destroying metabolic
timing coherence. Tumour acidification ($\Delta\mathrm{pH} \approx 0.3$) changes
the proton partition depth by $\Delta\Depth_{\ce{H+}} \approx 0.69$ (base-$e$)
and locks the cell into the glycolytic attractor. Late-stage cancer elevates
all three $\Sspace$-coordinates.

\medskip
\noindent\textbf{Neurodegeneration} ($\Delta\Sk > 0$): Ion pump failure raises
$\Sk$ by collapsing ionic partition depth differences. The cascade
(depolarisation $\to$ \ce{Ca^{2+}} influx $\to$ mitochondrial uncoupling
$\to$ apoptosis) is a positive-feedback partition depth collapse. In
Alzheimer's disease, amyloid-$\beta$ initiates the cascade through
Na$^+$/K$^+$-ATPase inhibition; in ALS, TDP-43 aggregation disrupts the
RNA partition gate, causing a mixed $\Sk$--$\Se$ deviation.

\medskip
\noindent\textbf{Aging} ($\Delta\Se > 0$): NPC degradation at rate
$k_{\mathrm{NPC}}(t) = k_0 e^{\alpha t}$ causes exponentially accelerating
nuclear-cytoplasmic mixing, epigenetic drift, and heterochromatin loss.
The $\Se$ trajectory follows Eq.~\eqref{eq:se_increase}, consistent with
methylation clock data.

\medskip
The framework predicts that $\Sspace$-deviation precedes clinical symptoms
in all three classes (Theorem~\ref{thm:sspace_predictive}), enabling
earlier diagnosis through ionic (for $\Sk$), metabolomic (for $\St$), and
epigenetic (for $\Se$) biomarkers.
\end{chapterbox}
